\section{Introduction}\label{sec:intro}

A zero-knowledge proof convinces a smart contract that \emph{some} statement
holds. It does not tell the contract what to do. Between the moment a verifier
returns \code{true} and the moment tokens move, the consuming contract decides
which public values to trust, which caller-supplied values to use, and which
state to update. If the action depends on a value the proof never fixed, a
perfectly sound proof authorizes an action nobody intended.

\para{A genuine proof, an unintended mint.}
Consider the \emph{ZKredential} compliance gateway from a 2025 hackathon
repository (Fig.~\ref{fig:zkredential}). The circuit proves that a public
wallet address satisfies a KYC predicate and exports an \code{isCompliant}
signal. The gateway checks the proof with an unmodified snarkjs-generated
Groth16 verifier, checks \code{isCompliant==1}, records the proof hash, and
then executes \code{mint(msg.sender, amount)}. Neither the recipient nor the
amount is a public input. We generated a real proof for one wallet, submitted
it from a different account, and chose the amount. On a local chain with
synthetic assets, the unchanged verifier and gateway accepted the proof and
minted 123\,456\,789\,012\,345\,678\,901\,234 units to the caller, while the account
named in the proof received
0~(\S\ref{sec:case-zkredential}).%
% src: reports/impact/REPORT.md#zkredential--propertyfy-gateway-mint (Effect row)
% src: reports/groundtruth/CLAIMS.md#A1
\
Every layer behaved correctly in isolation. The circuit is satisfiable only
for compliant witnesses; the verifier is the standard template; the contract
calls it. The flaw lives in the \emph{binding} between them.

\begin{figure}[t]
\centering
% F1: ZKredential flow. Facts: findings/zkredential/policy.md; DEV-F30 ComplianceGateway.sol:130-169
% (code lines abbreviated: n = pubSignals.length, x = pubSignals, pi = proof);
% effect values: reports/impact/REPORT.md#zkredential (observation.json); reports/groundtruth/CLAIMS.md#A1.
{\renewcommand{\code}[1]{\texttt{#1}}%
\begin{tikzpicture}[
  font=\scriptsize,
  box/.style={draw, rounded corners=2pt, align=left, inner sep=3pt, fill=white},
  bad/.style={draw=red!70!black, thick},
  arr/.style={-{Stealth[length=4pt]}, semithick},
  lab/.style={font=\tiny, align=center, fill=white, inner sep=1pt}
]
\node[box, text width=3.55cm] (circ) {\textbf{Circuit} (PropertyFy)\\
public: \code{walletAddress} $=W$\\
output: \code{isCompliant}\\
private: KYC and asset values};
\node[box, text width=3.2cm, right=0.45cm of circ] (ver) {\textbf{Verifier} (snarkjs, unchanged)\\
pairing check on $(\pi, x)$};
\node[box, text width=7.6cm, below=0.5cm of circ.south west, anchor=north west] (gw) {\textbf{Gateway} \code{verifyAndExecute(platform, $\pi$, $x$, MINT, amount)}\\[1pt]
\code{1 require(x.length == expected)}\\
\code{2 require(verifyProof($\pi$, $x$))} \hfill \textcolor{green!40!black}{fixes $W$, compliance}\\
\code{3 require(x[2] == 1)}\\
\code{4 used[keccak256($\pi$, $x$)] = true}\\
\code{5 mint(\textcolor{red!70!black}{msg.sender}, \textcolor{red!70!black}{amount})} \hfill \textcolor{red!70!black}{neither is a public input}};
\node[box, bad, text width=3.55cm, below=0.5cm of gw.south west, anchor=north west] (m) {\textbf{Caller} $M\neq W$, balance after:\\ \textbf{123456789012345678901234}};
\node[box, text width=3.2cm, right=0.45cm of m] (w) {\textbf{Proof wallet} $W$, balance after:\\ \textbf{0}};
\draw[arr] (circ) -- node[lab, above]{proof} (ver);
\draw[arr] (ver.south) -- node[lab, right]{accept} (ver.south |- gw.north);
\draw[arr] (gw.south -| m.north) -- node[lab, right]{mint} (m.north);
\end{tikzpicture}
}

\caption{ZKredential. The proof fixes the wallet $W$ and compliance; the
gateway mints a calldata \code{amount} to \code{msg.sender}~$=M\neq W$.
Neither effect parameter is a public input.}
\label{fig:zkredential}
\end{figure}

\para{Binding spans three artifacts.}
Whether a proof binds an action is a property of the triple (circuit,
generated verifier, consumer), not of any single file. The circuit fixes which
signals are public and how they are constrained; the generated verifier fixes
the arity, order and field-membership checks of the public vector; the
consumer fixes how that vector is assembled and which effects follow. A
recipient that is bound in the circuit is unbound if the consumer takes it from
calldata; a nullifier that is unique in the field is replayable if the verifier
accepts non-canonical encodings and the consumer keys a mapping on raw
\code{uint256} values. The SNARK SoK~\cite{sok-snarks} names this the
\emph{integration layer}, counts 13 of its 141 vulnerabilities there, and
reports that the layer has ``no specific defenses''.%
% src: RELATED-WORK.md#32-studies-benchmarks-and-adjacent-work (SoK row); reports/related/evidence/primary-text-greps.txt
\
A 2026 evaluation of six circuit tools on 70 Circom bugs explicitly leaves
integration to future work~\cite{zkp-tools-study}, and the largest executable
ZK bug collection, zkbugs, has 140 entries of which 0 are located in a
Solidity file~\cite{zkbugs}.%
% src: RELATED-WORK.md#32 (tools study row; zkbugs row); reports/related/evidence/zkbugs-census.txt
\
Contract analyzers cannot see the circuit; circuit analyzers cannot see the
consumer. The one published analyzer that targets ``ZKP-based smart contracts''
reads Solidity only~\cite{elsa}.%
% src: RELATED-WORK.md#31-tools (ELSA row: 4,086 .sol files, 0 circuit files)

\para{Naive cross-layer checks give false assurance.}
The obvious remedy is to compose existing checks. We built such compositions
and measured them. A consumer-and-verifier rule set with a same-policy
property-test harness certified 3 of 23 vulnerable synthetic cases as
safe~(\S\ref{sec:res-benchmark}).%
% src: reports/redteam/REPORT.md#synthetic-comparison (Arm A, vulnerable misses 3/23)
\
A circuit tool answered ``properly constrained'' on all 46 synthetic circuits,
vulnerable or not, because circuit uniqueness says nothing about what the
consumer does.%
% src: reports/redteam/REPORT.md#dev-corpus-and-native-tool-outcomes (Picus exit 8 for all 46)
\
On the historical zkDrops pair, whose vulnerable revision we reproduced as a
double redemption, both composed baselines emitted identical warnings on the
vulnerable and the repaired revision~(\S\ref{sec:res-pairs}).%
% src: paper/zk/data/pair-discrimination.json#pairs[0] (arm_A/arm_C identical_output)
\
Our own checker was not immune: earlier versions certified the vulnerable
zkDrops revision and a Semaphore revision with no verifier at all as
\VWS~(\S\ref{sec:design-failures}).%
% src: reports/groundtruth/REPORT.md#4-verdict-defects-observed-in-retained-dev-output

\para{This paper.}
We present a benchmark and measurement study of circuit--contract binding
flaws in Circom/Groth16 applications, and \tool{}, a cross-layer checker
designed around one rule: \emph{certify a binding only with path-sensitive,
position-exact evidence, and otherwise fail closed}. \tool{} reads the
circuit (source and compiled R1CS), the snarkjs-generated verifier and the
consuming contract, and returns, per binding property, \VIOL{},
\VWS{} or \UNK{} with a reason code. Every positive verdict names the evidence
it requires (Table~\ref{tab:evidence}); a missing artifact, an unresolved
modifier, or a guard that does not dominate the effect yields \UNK{}, never a
pass.

Our contributions are:
\begin{itemize}
\item \textbf{A taxonomy} of binding flaws grouped by what the proof fails to
  fix---the effect, its context, its uniqueness, the verifier and key
  configuration, and the interface shape---mapped onto the SoK's
  integration categories and extended with verifier-generation defects
  (\S\ref{sec:taxonomy}).
\item \textbf{A benchmark} that pins all three artifacts per case: 46 paired
  vulnerable/legitimate synthetic cases in 23 families, 56 development
  repositories including three historical vulnerable/fixed pairs, a 16-pair
  gate set, a preregistered 278-family held-out sample drawn from 1,384
  eligible families, and a passive census of 1,027 deployed verifiers
  (\S\ref{sec:benchmark}).%
% src: reports/redteam/REPORT.md#synthetic-comparison (23 vulnerable + 23 legitimate); reports/soundness/REPORT-2.md#full-dev-ledger (56 rows)
% src: reports/groundtruth/CLAIMS.md#A19,A20,A26; reports/groundtruth/REPORT.md#5-gate-manifest (16 rows)
\item \textbf{\tool{}}, a fail-closed checker with an explicit evidence
  contract per verdict, and a record of six distinct mechanisms by which
  certification went wrong without that contract---including one we found in
  the frozen version while preparing this paper (\S\ref{sec:design}).
\item \textbf{A measurement} of how single-layer tools, naive compositions and
  \tool{} behave on the benchmark: decided rates, warning restraint, unsafe
  certifications and pair discrimination (\S\ref{sec:results}), with
  held-out results \MISSING{held-out run, lane EVAL}.
\item \textbf{Validated findings}: two proof-to-effect violations reproduced
  end to end with real proofs and unchanged consumers on local chains, each
  with a repair that rejects the attack and preserves a legitimate run
  (\S\ref{sec:cases}).%
% src: reports/impact/REPORT.md (ZKredential and RacePass tables)
\end{itemize}

To our knowledge (bounded literature search, November 2026; method in
\S\ref{sec:related}), \tool{} is the first analyzer that takes a Circom
circuit's compiled constraints, its generated Groth16 verifier and the
consuming Solidity contract as joint input and reports, for each protected
effect, whether the accepted proof fixes it, failing closed to \UNK{} with a
reason code otherwise.%
% src: RELATED-WORK.md#6a (scoped priority sentence, hedges kept)
\
We make no claim that \tool{} detects more flaws than a purpose-built
contract-only rule: on our development set every \tool{} violation carries a
consumer- or verifier-level reason code (\S\ref{sec:res-ablation}). Its value,
as we show, lies in what it refuses to certify.
